\subsection{CASE OF GROUP GERMS}

Let \((\mathbf{G}, e, \theta, m)\) be a Lie group germ and \(\Omega\) the set of definition of m. Then \(\mathrm{T}(\Omega)\) is identified with an open subset of \(\mathrm{T}(\mathrm{G}) \times \mathrm{T}(\mathrm{G})\) and \(\mathrm{T}(m)\) is an analytic mapping of T(Ω) into T(G). It can be verified as in no. 2 that (T(G), 0e, T(θ), T(m)) is a Lie group germ. The products of G and T(G) are often written multiplicatively. The canonical projection of T(G) onto G is a morphism of Lie group germs. The restriction of Te(m) to Te(G) is the vector space addition of Te(G). The zero section of T(G) is an isomorphism of the Lie group germ G onto a Lie subgroup germ of T(G) which we identify with G. If f is a morphism of G into a Lie group germ H, T(f): T(G) → T(H) is a morphism of Lie group germs.

The mapping \(\phi\colon(g,u)\mapsto gu\) of \(G\times T_{e}(G)\) into \(T(G)\) is an isomorphism of the trivial vector bundle \(G\times T_{e}(G)\) with base space \(G\) onto the vector bundle \(T(G)\) ; for \(\phi\) and \(\phi^{-1}\) are analytic and are vector bundle morphisms, so that it suffices to apply Differentiable and Analytic Manifolds, R, 7.2.1. (The proof of no. 2 could also be adapted.) The isomorphism \(\phi^{-1}\) is called the left trivialization of \(T(G)\) . The inverse isomorphism of the mapping \((g,u)\mapsto ug\) is called the right trivialization.

Let X be a manifold of class \(C^{r}\) and \(\psi\) a law chunk of left operation of class \(C^{r}\) of G on X. Then \(\mathrm{T}(\psi)\) is a law chunk of left operation of class \(C^{r-1}\) of \(\mathrm{T}(\mathrm{G})\) on T(X) extending \(\psi\) . Formulae (7), (8) and (9) remain valid if gx is defined. If I is an open subset of K containing 0, if \(\gamma: I \to G\) is a mapping of class \(C^{r}\) such that \(\gamma(0) = e\) and if \(a = \mathrm{T}_{0}(\gamma)1\) , the vector field \(x \mapsto ax\) defined on X is the vector field defined by the mapping \((\lambda, x) \mapsto \gamma(\lambda)x\) in the sense of Differentiable and Analytic Manifolds, R, 8.4.5.

